\section{Spectral functions: heat kernel, zeta function, resolvent}
\label{sec2}
\subsection{Differential geometry and the operators of Laplace type}
\label{s2geo} 
Let $M$ be a smooth compact Riemannian manifold\footnote{This simply
means that we assume that there is a positive definite metric
tensor $g_{\mu\nu}$ on $M$.}
of dimension $n$ with smooth boundary $\partial M$.
We shall also consider the case when the boundary
$\partial M$ is empty. Let $V$ be a vector 
bundle\footnote{Many physicists strongly dislike vector bundles.
Nevertheless, there are two good reasons for using the fibre
bundles in this paper (in parallel with more familiar notations of matrix-valued
functions, gauge fields, etc). 
First, our simplifying comments and
examples may help the reader to understand mathematical
literature on the subject. Second, one of the main ideas of this
report is to reveal some universal structures behind the heat
kernel expansion. In particular, we shall see that there is not
much difference between different spins and symmetry groups.
The vector bundle language seems to be the most adequate language
to achieve this goal. The reader may consult the excellent review paper 
by Eguchi, Gilkey and Hanson \cite{Eguchi:1980jx}.
} over $M$. This means that there is a vector space
attached to each point of the manifold. For example, this
could be a representation space of a gauge group or of the
space-time symmetry group.
Sections of $V$ are smooth functions bearing
a discrete index which could correspond to internal
or spin degrees of freedom. 

We study differential
operators on $V$. We restrict ourselves to second order
operators of the Laplace type. Locally such operators can
be represented as:
\begin{equation}
D=-(g^{\mu\nu}\partial_\mu\partial_\nu 
+a^\sigma\partial_\sigma +b) \label{D1}
\end{equation}
where $g^{\mu\nu}$ is the inverse metric tensor on $M$;
$a^\sigma$ and $b$ are matrix valued functions on $M$.
There is a unique connection on $V$ and a unique 
endomorphism $E$ of $V$ (another matrix valued function)
so that
\begin{equation}D=-(g^{\mu\nu}\nabla_\mu\nabla_\nu +E)\, ,
\label{laplaceb}\end{equation}
where the covariant derivative $\nabla =\nabla^{[R]} +\omega$
contains both Riemann $\nabla^{[R]}$ and ``gauge'' (bundle) $\omega$
parts.
We may express:
\begin{eqnarray}
    &&\omega_\delta=\frac 12 g_{\nu\delta}(a^\nu+g^{\mu\sigma}\Gamma_{\mu
    \sigma}{}^{\nu }I_{V}) \quad{\rm and}\label{oab}\\
    &&E=b-g^{\nu \mu}(
    \partial_\mu\omega_\nu+\omega_\nu\omega_\mu
     -\omega_\sigma\Gamma_{\nu\mu}{}^\sigma)\,, \label{Eab}
\end{eqnarray}
where
\begin{equation}
\Gamma_{\nu\mu}{}^\sigma =g^{\sigma\rho} \frac 12 \left(
\partial_\mu g_{\nu\rho}+\partial_\nu g_{\mu\rho} -
\partial_\rho g_{\mu\nu} \right) \label{Christ}
\end{equation}
is the Christoffel symbol. 
$I_{V}$ is the unit operator on $V$. 

It is important to understand how general the notion of the Laplace
type operator is. The most obvious restriction on $D$ is that it is
of second order, i.e. it contains second derivatives, but does not
contain higher derivative parts. In this paper we shall also consider
first order operators (Dirac operators, for example). Second, the
operator $D$ is a {\it partial differential operator}. This excludes
negative or fractional powers of the derivatives. The operators containing
such structures are called {\it pseudo-differential operators}.
More information on spectral theory for these operators can be found
in \cite{Seeley:1968,Grubb:1996,Grubb:2002}. 
Third, the operator (\ref{D1}) has a
scalar principal part. This means that the second derivatives in
(\ref{D1}) are contracted with the metric, and the internal index structure
of the second derivative term is trivial. Such operators are also
called {\it minimal}. Non-minimal operators will be briefly
considered in sec\ \ref{s4nonmin}.

We can define local invariants associated with $\omega$
and $g$. Let 
\begin{equation}
{R^\mu}_{\nu\rho\sigma}=\partial_\sigma \Gamma_{\nu\rho}^\mu -
\partial_\rho \Gamma_{\nu\sigma}^\mu +
\Gamma_{\nu\rho}^\lambda \Gamma_{\lambda\sigma}^\mu -
\Gamma_{\nu\sigma}^\lambda \Gamma_{\lambda\rho}^\mu 
\label{Riemann}
\end{equation}
be the Riemann curvature tensor, let $R_{\mu\nu}:={R^\sigma}_{\mu\nu\sigma}$
be the Ricci tensor, and let $R:=R^\mu_\mu$ be the scalar
curvature. 
With our sign
conventions, $R =2$ on the unit sphere $S^2$ in the Euclidean
space. 
Let
Roman indices $i$, $j$, $k$, and $l$ range from 1 through
the dimension $n$ of the manifold and index a local orthonormal frame
(vielbein) $\{ e_1,...,e_n\}$ for the tangent space of the manifold.
In components we have: $e_j^\mu e_k^\nu g_{\mu\nu}=\delta_{jk}$,
$e_j^\mu e_k^\nu \delta^{jk}=g^{\mu\nu}$. The inverse vielbein
is defined by the relation $e_\mu^j e^\mu_k=\delta_k^j$. These two
objects, $e_\mu^k$ and $e^\nu_j$, will be used to transform 
``curved'' indices ($\mu$, $\nu$, $\rho$) to ``flat'' ones
($i$, $j$, $k$) and back. In Euclidean space there is no distinction
between upper and lower flat indices.

As usual, the Riemann part of the covariant derivatives contains the
Christoffel connection so that
\begin{equation}
\nabla_\mu^{[R]}v_\nu = \partial_\mu v_\nu -\Gamma_{\mu\nu}^\rho v_\rho
\label{riemcov}
\end{equation}
for an arbitrary vector $v_\nu$. To extend this derivative to the objects
with flat indices one has to use the spin-connection $\sigma_\mu$:
\begin{equation}
\nabla_\mu v^j = \partial_\mu v^j +\sigma_\mu^{jk}v_k \,.\label{spinder}
\end{equation}
The condition $\nabla_\mu e_\nu^k=0$ yields
\begin{equation}
\sigma_\mu^{kl}=e_l^\nu \Gamma_{\mu\nu}^\rho e_\rho^k -
e^\nu_l \partial_\mu e_\nu^k \,.\label{spincon}
\end{equation} 

Let $\Omega_{\mu\nu}$ be the field strength of the connection
$\omega$:
\begin{equation}
\Omega_{\mu\nu}=\partial_\mu\omega_\nu - \partial_\nu\omega_\mu
+\omega_\mu\omega_\nu -\omega_\nu\omega_\mu \, .
\label{Omega}
\end{equation}
The covariant derivative $\nabla$ acts on both space-time and internal
indices. For example,
\begin{equation}
\nabla_\rho \Omega_{\mu\nu}=\partial_\rho \Omega_{\mu\nu}
-\Gamma_{\rho\mu}^\sigma \Omega_{\sigma\nu}
-\Gamma_{\rho\nu}^\sigma \Omega_{\mu\sigma}
+[\omega_\rho ,\Omega_{\mu\nu}] \,.
\end{equation}

If the boundary $\partial M$ is non-empty we have more invariants.
Let $e_n$ be inward pointing unit vector field. Let Roman indices $a$, $b$,
$c$ and $d$ range from 1 to $n-1$ and index a local orthonormal frame
for the tangent bundle of $\partial M$. Let $L_{ab}:=\Gamma^n_{ab}$
be the second fundamental form (extrinsic curvature) of the boundary.
We use the Levi-Civita (spin) connections and
the connection $\omega$ to covariantly differentiate tensors of all types. Let
\lq ;\rq\ denote multiple covariant differentiation with respect to 
the Levi-Civita connection of $M$ and let
\lq:\rq\ denote multiple tangential covariant differentiation 
on the boundary with respect to the
Levi-Civita connection of the boundary; the difference
 between `;' and `:' is measured by the second
fundamental form. Thus, for example, $E_{;a}=E_{:a}$ since 
there are no tangential indices in $E$ to be
differentiated. On the other hand, 
$E_{;ab}\ne E_{:ab}$ since the index $a$ is also being
differentiated. More precisely $E_{;ab}=E_{:ab}-L_{ab}E_{;n}$.
Since $L$ is only defined on the boundary, this
tensor can only be differentiated tangentially.

Consider an example of the circle $S^1$ in the plane $\mathbb{R}^2$.
The line element has the form
\begin{equation}
(ds)^2 = (dr)^2 +r^2 (d\theta)^2 \,,\label{lineS2}
\end{equation}
where $0\le r<\infty$, $0\le \theta <2\pi$. Then $g_{\mu\nu}=
{\rm diag} (1,r^2)$. $S^1$ is defined by the condition $r=r_0$.
We may choose $e_\theta^1=r$, $e_r^2=-1$
(the minus sign appears because $e^2$ is an {\it inward}
pointing unit vector for the disk with the boundary $S^1$).
Then the second fundamental form of $S^1$ is
\begin{equation}
L_{11}=\left.
e_r^2 e_1^\theta e_1^\theta \Gamma_{\theta\theta}^r\right|_{r=r_0}=
\frac 1{r_0} \,.\label{LonS1}
\end{equation}
In general, on $S^{n-1}$ considered as a boundary of the ball in
$\mathbb{R}^n$ the extrinsic curvature is $L_{ab}=\frac 1{r_0} \delta_{ab}$.

If the boundary $\partial M$ is non-empty, one has to
define boundary conditions for the field $\phi$. A convenient
way to write them down is:
\begin{equation}
\mathcal{B}\,\phi =0 \label{boucon}
\end{equation}
where $\mathcal{B}$ is called the boundary operator. 
In general, the operator $\mathcal{B}$ calculates a linear combination
of the boundary data for any given function $\phi$. If $D$ is of Laplace
type, the boundary data include value of the function at the
boundary and value of it's first normal derivative.
The most frequently used
choices are the Dirichlet and Neumann boundary operators
which we denote $\mathcal{B}^-$ and $\mathcal{B}^+$ respectively:
\begin{eqnarray}
&&\mathcal{B}^-\phi =\phi\vert_{\partial M} \, ,\label{Dirbc} \\
&&\mathcal{B}^+\phi =(\phi_{;n}+S\phi )\vert_{\partial M} \, ,
\label{Neubc}
\end{eqnarray}
where $S$ is a matrix valued function defined on $\partial M$.
The boundary conditions (\ref{Neubc}) are also
called Robin or generalised Neumann. In some literature they
are called mixed boundary conditions. We shall not use this latter
terminology. The name ``mixed'' is reserved for another type
of boundary conditions. Let $\Pi^-$ and $\Pi^+$ be two
complementary projectors defined on $V\oB$, $(\Pi^\pm)^2=\Pi^\pm$,
$\Pi^++\Pi^-=I$.
There is a decomposition
$ V|_{\partial M}=V_N\oplus V_D$, where $V_{N,D}
=\Pi^\pm V|_{\partial M}$. Decompose also 
$\phi=\phi^N\oplus\phi^D$ and set
\begin{equation}
\mathcal{B}\phi:=\phi^D\oplus(\phi^N_{;n}+S\phi^N)
\oB \, .
\label{mixbc}
\end{equation}
The matrix valued function (endomorphism) $S$ acts on $V_N$ only, 
$S=\Pi^+S=S\Pi^+$. In other words, we define Dirichlet boundary
conditions on $V_D$ and generalised Neumann boundary conditions on $V_N$.
For obvious reason the boundary conditions (\ref{mixbc})
will be called mixed. In  sec. \ref{sec3} we shall see that natural
boundary conditions for spinor and vector fields are of this type.
%%%%%%%%%%
\subsection{Spectral functions}\label{s2spec}
For the boundary conditions considered in this section as well
as on manifolds without a boundary the operator 
$\exp (-tD)$ with positive $t$ is trace class on the space
of square integrable functions $L^2(V)$. This means that for
an auxiliary smooth function $f$ on $M$ 
\begin{equation}
K(t,f,D)={\rm Tr}_{L^2}(f\exp (-tD)) 
\label{heatTr}
\end{equation}
is well defined. We also write
\begin{equation}
K(t,f,D)=\int_M d^nx \sqrt{g} \ptr K(t;x,x;D) f(x) \,,\label{KxxD}
\end{equation}
where $K(t;x,x;D)$ is an $y\to x$ limit of the fundamental
solution
$K(t;x,y;D)$ of the heat equation (\ref{heateq})
with the initial condition (\ref{inicon}).
If there is a boundary, the kernel $K(t;x,y;D)$ should also
satisfy some boundary conditions $\mathcal{B}_xK(t;x,y;D)=0$
in one of the arguments. We stress that $K(t;x,y;D)$ is a matrix
in the internal indices. $\ptr$ denotes the trace over these indices.

Let $D$ be self-adjoint. This implies that in a suitable basis in the
internal space the matrix $\omega_\mu$ is anti-hermitian and $E$ is
hermitian. Let $\{ \phi_\lambda \}$ be a complete basis of orthonormal
eigenfunctions of the operator $D$ corresponding to the eigenvalues
$\{\lambda\}$. Then
\begin{equation}
K(t;x,y;D)=\sum_\lambda \phi^\dag_\lambda (x)\phi_\lambda (y)
e^{-t\lambda} \,.\label{hkeigen}
\end{equation}

We shall almost exclusively work either on manifolds without boundaries, or
on manifolds with boundaries with the fields subject to local boundary
conditions (\ref{Dirbc}), (\ref{Neubc}), or (\ref{mixbc}). In all these
cases there is
an asymptotic expansion as $t\downarrow 0$ 
\cite{Seeley:1966,Seeley:1969,Greiner:1971}\footnote{These papers contain
also a method allowing to calculate the coefficients of the expansion.
The method is, however, too complicated to use it on practice.
}:
\begin{equation}
{\rm Tr}_{L^2} (f\exp (-tD)) 
\cong\sum_{k\ge0}t^{(k-n)/2}a_k(f,D)\,.
\label{asymptotex}
\end{equation}
This expansion is valid for almost all boundary conditions appearing
in applications to physics. There are, however, some exceptions which
will be discussed in sec.\ \ref{s5other}.

The coefficients $a_k$ and the coefficients $b_k$ introduced in the
previous section (cf. (\ref{sHKexp})) are related by the equation
($m=0$ for the simplicity):
\begin{equation}
a_k(f,D)=(4\pi )^{-n/2} \iM b_k(x,x)f(x) \,.\label{akbk}
\end{equation}
Note, that the definition (\ref{sHKexp}) is valid on flat manifolds
without boundary only (although generalisations to other cases
are possible).

The key property of the heat kernel coefficients $a_k$ is that
they are locally computable in most of the cases. This means that
they can be expressed in term of the volume and boundary integrals
of local invariants. 

For a positive operator $D$ one can define the zeta function
\cite{Mina,Seeley:1968}
by the equation:
\begin{equation}
\zeta (s,f,D)= {\rm Tr}_{L^2} (f\, D^{-s}) \,.\label{defzeta}
\end{equation}
The zeta function is related to the heat kernel by the integral
transformation
\begin{equation}
\zeta (s,f,D)=\Gamma (s)^{-1}\int\limits_0^\infty dt\, t^{s-1}
K(t,f,D) \,.\label{zetaheat}
\end{equation}
This relation can be inverted,
\begin{equation}
K(t,f,D)=\frac 1{2\pi i}\oint ds\, t^{-s}\Gamma(s)\zeta (s,f,D)\,,
\label{heatzeta}
\end{equation}
where the integration contour encircles all poles of the integrand.
Residues at the poles can be related to the heat kernel coefficients:
\begin{equation}
a_k(f,D)={\rm Res}_{s=(n-k)/2} \left( \Gamma(s)\zeta (s,f,D) \right)\,.
\label{aRes}
\end{equation}
In particular,
\begin{equation}
a_n(f,D)= \zeta (0,f,D) \,.\label{anz0}
\end{equation}

Zeta-functions can be used to regularize the effective action
\cite{Dowker:1976tf,Hawking:1977ja}. The regularization is achieved
by shifting the power of $t$ in (\ref{WKtD}):
\begin{equation}
W_s=-\frac 12 \tilde\mu^{2s}\int_0^\infty \frac {dt}{t^{1-s}} 
K(t,D) \,,\label{Wzreg}
\end{equation}
where $\tilde\mu$ is a constant of the dimension of mass introduced to keep
proper dimension of the effective action. The regularization is removed
in the limit $s\to 0$. Eq.\ (\ref{Wzreg}) can be considered as a
definition of the regularized effective action without any reference to (\ref{WKtD}).
One can also rewrite (\ref{Wzreg}) in terms of the zeta function:
\begin{equation}
W_s=-\frac 12 \tilde\mu^{2s} \Gamma(s)\zeta (s,D) \,,
\label{WzGam}
\end{equation}  
where $\zeta (s,D):=\zeta (s,1,D)$.

The gamma function has a simple pole at $s=0$:
\begin{equation}
\Gamma (s)=\frac 1s -\gamma_E +\mathcal{O}(s)\,,\label{Gammas}
\end{equation}
where $\gamma_E$ is the Euler constant. The regularized effective action
(\ref{WzGam}) has also a pole at $s=0$:
\begin{equation}
W_s=-\frac 12 \left( \frac 1s -\gamma_E +\ln \tilde\mu^2 \right)
\zeta (0,D) - \frac 12 \zeta'(0,D) \,.\label{Wzpole}
\end{equation}
According to (\ref{anz0}) the divergent term in the zeta function 
regularization is proportional to $a_n(D)$ (cf. (\ref{divWLam}) for another
regularization scheme). The pole term in (\ref{Wzpole}) has to be removed
by the renormalization. The remaining part of $W_s$ at $s=0$ is the 
renormalised effective action:
\begin{equation}
W^{\rm ren}= -\frac 12 \zeta'(0,D)-\frac 12 \ln (\mu^2 )
\zeta (0,D) \,,\label{Wzren}
\end{equation}
where we have introduced a rescaled parameter 
$\mu^2=e^{-\gamma_E}\tilde\mu^2$. In this approach $\mu^2$ describes the
renormalization ambiguity which must be fixed by a suitable
normalisation condition. Let us remind that here we are working on a
compact manifold. On non-compact manifolds $\zeta (s,D)$ may have divergent
contributions proportional to the volume. Such divergences are usually
removed by the subtraction of a ``reference'' heat kernel (see sec.\
\ref{s6sinpot}).

Together with (\ref{defW}) equation (\ref{Wzren}) yields a definition
of the functional determinant for a positive elliptic second order
operator which is frequently used in mathematics:
\begin{equation}
\ln \det (D)=-\zeta'(0,D) -\ln (\mu^2) \zeta (0,D) \,.\label{detzeta}
\end{equation}

Note, that the definitions (\ref{defzeta}), (\ref{detzeta})
are valid for positive operators
only. Elliptic 2nd order differential operators have at most finite number
of zero and negative modes\footnote{On manifolds with boundaries one has to
require {\it strong ellipticity} of the boundary value problem (cf.\
sec.\ \ref{s5other}).} which must be treated separately. However, one 
can extend the definition of the zeta function to operators with negative
modes:
\begin{equation}
\zeta (s,D)=\sum |\lambda |^{-s} \,,\label{zetaneg}
\end{equation}
where the sum extends over all non-zero eigenvalues $\lambda$.
One can also define another spectral function in a similar way:
\begin{equation}
\eta (s,D)=\sum {\rm sign}\, (\lambda ) |\lambda |^{-s} \,.\label{etadef}
\end{equation}
This function is especially useful in spectral theory of Dirac type
operators where $\eta (0,\Dir )$ measures asymmetry of the spectrum.

Another function which is frequently used especially in the mathematical
literature is the {\it resolvent} (or, more precisely, its' powers):
\begin{equation}
{\rm R}^l(z):=(D+z^2)^{-l} \,.\label{resolvent}
\end{equation}
If $D$ is on operator of Laplace type subject to ``good''
boundary conditions\footnote{A more precise meaning of this 
restriction will be explained in sec. \ref{s5other}.}
 and if $l$ sufficiently large ($l>n/2$)
there is a full asymptotic expansion
\begin{equation}
{\rm Tr} ({\rm R}^l(z))=\sum_k 
\frac{\Gamma \left( l+\frac{k-n}2 \right)}{\Gamma (l)}
a_k(D) z^{-2l+n-k} \label{asresol}
\end{equation} 
as $z\to \infty$. The coefficients $a_k$ are the same as in the
heat kernel expansion (\ref{asymptotex}).
%%%%%%
\subsection{Lorentzian signature}\label{s2Lorentz}
Locality of the heat kernel coefficients in the Euclidean domain
can be easily understood by examining the free heat kernel
(\ref{simplestHK}). For small $t$ the first term in the exponential
strongly suppresses non-local contributions. For Lorentzian 
metrics the squared distance function $(x-y)^2$ is no longer
positive definite. Therefore, the simple arguments given above
do not work. A partial solution to this problem is to consider
the ``Schr\"{o}dinger'' equation
\begin{equation}
\left( i\partial_\tau -D \right) \tilde K(\tau;x,y;D)=0
\label{Schr}
\end{equation}
for the kernel $\tilde K$ instead of the heat conduction
equation (\ref{heateq}). Then $\tilde K$ oscillates at large
distances. However, even though non-local contributions
to $\tilde K$ oscillate furiously as $\tau\to 0$ they are
not small. Consequently, local asymptotic series do not
exist in many cases. A discussion on this point
can be found in chapter 9 of \cite{Fulling:1989nb}.

Of course, the heat kernel expansion can be used also
on Lorentzian manifolds at least for the renormalization
theory where non-local terms are of less importance.
Counterterms are still defined by the same heat kernel
coefficients with the same functional dependence on
the Lorentzian metric. Explicit definitions with the imaginary
``proper time'' $\tau$ can be found in 
\cite{DeWitt:1965,Birrell:1982ix}. One should note that some
background fields receive an imaginary phase when being
continued to the Euclidean domain (cf. secs.\ \ref{s3string}
and \ref{s3spin}).

